\begin{helper}[Campanato decay for the gradient from a comparison bundle]
Let $n\ge2$, $p>1$, $0<\alpha\le1$, $\delta'>0$ and
\[
0<\tau<p\alpha\theta,\qquad \theta=\frac{\delta'}{n+p\alpha} .
\]
Then there are $\lambda_0\in(0,p)$ and $R_2>0$ with $R_2\le1$, depending only on
$n,p,\alpha,\delta'$ and $\tau$, with the following property.  (The upper bound
on $\lambda_0$ is \emph{strict}; this is what the consumers need, because
Campanato's theorem (C5) is available at the exponent $\lambda=n+\lambda_0$ only
for $\lambda<n+p$.  The paper's Theorems 4 and 6 state $0<\lambda_0\le p$, which
is weaker and which they keep.)

Let $c_1,c_3>0$.  Then there is $c_{10}>0$, depending only on
$n,p,\alpha,\delta',\tau,c_1$ and $c_3$, such that the following holds.

Let $0\le\varepsilon\le1$, let $\Omega\subseteq\R^n$ be open, let $\Omega'$ be
open with $\overline{\Omega'}\subseteq\Omega$, let $R_1>0$ and let
$u\in W^{1,p}(\Omega)$ with weak gradient $g=\nabla u$ (\noderef{MemW1pOn}).
Write
\[
\Lambda(x_0,R)=R^{-n}\int_{B_R(x_0)}(1+|\nabla u|^2)^{\frac p2}dx .
\]
Assume
\begin{align}
\Lambda(x_0,R)&\le c_1R^{-\tau}\qquad (x_0\in\Omega',\ 0<R\le R_1),\tag{4.15}
\end{align}
and that for every $x_0\in\Omega'$ and every $0<R\le R_1$ with
$\overline B_R(x_0)\subseteq\Omega$ there is a comparison map
$v\in W^{1,p}\bigl(B_R(x_0)\bigr)$ with weak gradient $g_v$
(\noderef{MemW1pOn}) such that
\begin{align}
\int_{B_\rho(x_0)}\bigl|g_v-(g_v)_{B_\rho(x_0)}\bigr|^p dx
&\le c_3\Bigl(\frac\rho R\Bigr)^{p\alpha}\rho^{\,n}\Lambda(x_0,R)
\qquad(0<\rho\le R/4),\tag{4.16}\\
\int_{B_R(x_0)}|\nabla u-g_v|^p\,dx&\le R^{\delta'+n}\Lambda(x_0,R).\tag{4.17}
\end{align}
Then for every $x_0\in\Omega'$, every $R$ with $0<R\le\min\{R_1,R_2\}$ and
$\overline B_R(x_0)\subseteq\Omega$, and every $\rho>0$ with
$2\rho\le R^{1+\theta}$,
\[
\int_{B_\rho(x_0)}\bigl|\nabla u-(\nabla u)_{B_\rho(x_0)}\bigr|^p dx
\le c_{10}\,\rho^{\,n+\lambda_0}.
\tag{4.19}
\]
Here $(h)_{B_\rho(x_0)}=|B_\rho|^{-1}\int_{B_\rho(x_0)}h\,dx$ is the vector mean.

\noindent\textbf{On the quantifier order.}  $\lambda_0$ and $R_2$ are bound
before $c_1$ and $c_3$, and only $c_{10}$ is allowed to depend on them.  This is
the dependence the proof actually produces --- $\lambda_0=\eta/(1+\theta)$ and
$R_2=\min\{1,2^{-1/\theta}\}$ are built from $n,p,\alpha,\delta',\tau$ alone ---
and it is what the consumers need: in Theorem~4's and Theorem~6's route the
Morrey constant is produced by Theorem~3 at a stage where the decay exponent and
the radius threshold have already been fixed, so a $\lambda_0$ depending on
$c_1$ would be useless to them.

\noindent\textbf{On the membership $v\in W^{1,p}(B_R(x_0))$.}  The comparison
map of \S4 and \S5 is a $W^{1,p}$ solution of a frozen Dirichlet problem, so
this is what every producer of the bundle actually has; the thesis writes
(4.16) of the thesis about $\nabla v$ and (4.17) of the thesis about $\nabla w$ with
$w=u-v$ (thesis p.~43--44).  It is recorded explicitly because the proof below
\emph{adds} the two integrals bounded by (4.16) of the thesis and (4.17) of the thesis,
and an inequality bounding an integral does not by itself make the integrand
integrable.

\noindent\textbf{On the range $\rho\le R/4$ in (4.16) of the thesis.}  The thesis
states (4.16) of the thesis for $\rho\le R/2$ (p.~44).  It is assumed here on the
smaller range $\rho\le R/4$, which makes this node's hypothesis \emph{weaker}
and so its statement stronger, for two reasons.  First, the proof below uses
(4.16) of the thesis at one pair $(\rho,R')$ only, namely
$\rho=\frac12R'^{\,1+\theta}$, and the choice of $R_2$ made in the proof puts
that single $\rho$ below $R'/4$; so nothing on the range $(R/4,R/2]$ is ever
consumed, by either range convention.  Second, the producers of the bundle in
\S4 and \S5 deliver exactly the $R/4$ form: they obtain the oscillation
estimate from Theorem~2 on $B_{R/4}(x_0)$, because the sup bound that
Theorem~2 consumes is itself produced by Theorem~1 on a half-ball.  Assuming
the $R/4$ form therefore costs its consumers nothing and saves them a radius
conversion.

\noindent\textbf{On the range of $\rho$ in the conclusion.}  The conclusion is
stated for the radii $\rho$ with $2\rho\le R^{1+\theta}$ at an admissible
comparison radius $R$, and not for all $\rho\le\min\{R_1,R_2\}$, because that is
exactly the range the hypotheses reach.  The thesis derives (4.19) of the thesis by
putting $\rho=\frac12R^{1+\theta}$ (p.~45), so the estimate at a given $\rho$
consumes the comparison bundle at the \emph{larger} radius
$R=(2\rho)^{1/(1+\theta)}$ --- larger because $\theta>0$ and $2\rho\le1$ --- and
the bundle is available only for radii $\le R_1$ whose closed ball lies in
$\Omega$.  A conclusion quantified over all $\rho\le\min\{R_1,R_2\}$ would
assert the decay at radii whose associated comparison radius exceeds $R_1$,
where the hypotheses say nothing.  Conversely nothing is lost: as $R'$ runs over
$(0,R]$ the number $\frac12R'^{\,1+\theta}$ runs over all of
$(0,\frac12R^{1+\theta}]$, so every $\rho$ in the stated range is of the thesis'
special form for exactly one $R'\le R$, and the containment
$\overline B_R(x_0)\subseteq\Omega$ passes to $\overline B_{R'}(x_0)$.  A
consumer fixes $R$ from the geometry of its subdomain once and then has
(4.19) of the thesis for all small $\rho$, which is what Campanato's
characterisation of the H\"older classes consumes.
\end{helper}
\begin{proof}
Throughout, $B_t=B_t(x_0)$, $|E|$ is Lebesgue measure, $\omega_n=|B_1(0)|$, and
all non-integer powers are real powers of positive numbers.  Put
\[
\theta=\frac{\delta'}{n+p\alpha}>0,
\qquad
\eta=p\alpha\theta-\tau>0,
\qquad
\lambda_0=\frac\eta{1+\theta},
\]
so that $n\theta+p\alpha\theta=\delta'$ by the definition of $\theta$, and
$\eta>0$ is exactly the hypothesis $\tau<p\alpha\theta$; hence $0<\lambda_0$.
Also $\lambda_0<p$: since $\tau>0$ we have $\eta<p\alpha\theta$, and
$\alpha\le1$ with $\theta>0$ gives $\alpha\theta\le\theta<1+\theta$, so
\[
\lambda_0=\frac\eta{1+\theta}<\frac{p\alpha\theta}{1+\theta}<p ,
\]
dividing the inequality $p\alpha\theta<p(1+\theta)$ by $1+\theta>0$.  Put also
\[
R_2=\min\bigl\{1,\ 2^{-1/\theta}\bigr\},
\]
so that $0<R_2\le1$, since $2^{-1/\theta}>0$.  Both $\lambda_0$ and $R_2$ are
built from $n,p,\alpha,\delta'$ and $\tau$ alone --- $\theta$ and $\eta$ are,
and $\lambda_0$ and $R_2$ are functions of those --- so they may be, and are,
fixed before $c_1$ and $c_3$, which is the quantifier order the statement
prescribes.

Now let $c_1,c_3>0$ be given and put
\[
c_{10}=2^{p}\cdot2^{p-1}c_1\bigl(c_32^{-p\alpha-n}+1\bigr)\cdot2^{\,n+\lambda_0}>0 ,
\]
which depends only on $n,p,\alpha,\delta',\tau,c_1,c_3$ and is fixed before
$\varepsilon,\Omega,\Omega',u,g,R_1$.

\emph{Step 0: an elementary inequality.}  Since $p\ge1$, convexity of
$\lambda\mapsto\lambda^{p}$ on $[0,\infty)$ gives
$\bigl(\frac{X+Y}2\bigr)^{p}\le\frac12(X^{p}+Y^{p})$, i.e.
\[
(X+Y)^{p}\le2^{p-1}\bigl(X^{p}+Y^{p}\bigr)\qquad(X,Y\ge0),
\]
and therefore, by the triangle inequality,
\[
|a+b|^{p}\le2^{p-1}\bigl(|a|^{p}+|b|^{p}\bigr)\qquad(a,b\in\R^n).
\tag{E1}
\]

\emph{Step 1: the comparison radius attached to $\rho$.}  Let $x_0\in\Omega'$,
let $0<R\le\min\{R_1,R_2\}$ with $\overline B_R\subseteq\Omega$, and let
$\rho>0$ with $2\rho\le R^{1+\theta}$.  Put
\[
R'=(2\rho)^{\frac1{1+\theta}},\qquad\text{so that}\qquad
R'^{\,1+\theta}=2\rho,\qquad \rho=\tfrac12R'^{\,1+\theta}.
\]
Then $R'>0$, and $R'\le R$ because $2\rho\le R^{1+\theta}$ and
$\lambda\mapsto\lambda^{1/(1+\theta)}$ is nondecreasing on $[0,\infty)$, the
exponent $1/(1+\theta)$ being positive, and because
$(R^{1+\theta})^{1/(1+\theta)}=R$.  Hence
\[
0<R'\le R\le\min\{R_1,R_2\}\le R_2\le1,
\qquad
\overline B_{R'}\subseteq\overline B_R\subseteq\Omega .
\]
Three consequences will be used.  First $R'\le R_1$, so (4.15) of the thesis is
available at the radius $R'$ and the comparison bundle exists at $R'$: there is
$v\in W^{1,p}(B_{R'})$ with weak gradient $g_v$ satisfying (4.16) of the thesis and
(4.17) of the thesis on $B_{R'}$.  Second $R'\le R_2\le2^{-1/\theta}$, so, raising to
the power $\theta>0$,
\[
R'^{\,\theta}\le\bigl(2^{-1/\theta}\bigr)^{\theta}=2^{-1}=\tfrac12 ,
\]
whence, from $R'^{\,1+\theta}=R'\cdot R'^{\,\theta}$,
\[
\rho=\tfrac12R'\cdot R'^{\,\theta}\le\tfrac12R'\cdot\tfrac12=\tfrac14R' ,
\]
so (4.16) of the thesis may be applied at the pair $(\rho,R')$, its range
$0<\rho\le R'/4$ being met.  Third, $\rho\le R'/4\le R'$, so
$B_\rho\subseteq B_{R'}$.

This is the step the thesis leaves implicit, and it is where the ``radii not of
the special form'' difficulty dissolves: the thesis fixes $R$ and produces the
estimate at the single radius $\frac12R^{1+\theta}$, but the map
$R'\mapsto\frac12R'^{\,1+\theta}$ is a continuous increasing bijection of
$(0,R]$ onto $(0,\frac12R^{1+\theta}]$, so solving for $R'$ instead of for
$\rho$ covers \emph{every} admissible $\rho$ exactly, with no interpolation
between special radii and hence no appeal to monotonicity in $\rho$.

\emph{Step 2: the estimate (4.18) at the pair $(\rho,R')$.}  Put
$c=(g_v)_{B_\rho}$ and $\Lambda=\Lambda(x_0,R')$.  By (E1) with
$a=g_v-c$ and $b=g-g_v$, i.e.\ $g-c=(g_v-c)+(g-g_v)$,
\[
|g-c|^{p}\le2^{p-1}\bigl(|g_v-c|^{p}+|g-g_v|^{p}\bigr)
\]
pointwise; integrating over $B_\rho$ and enlarging the domain of the second term
to $B_{R'}$ (the integrand is nonnegative and $B_\rho\subseteq B_{R'}$),
\[
\int_{B_\rho}|g-c|^{p}dx
\ \le\ 2^{p-1}\int_{B_\rho}|g_v-c|^{p}dx+2^{p-1}\int_{B_{R'}}|g-g_v|^{p}dx .
\]
Every integral here is finite, and each from a \emph{membership} rather than
from a bound: $|g|^{p}$ is integrable on $\Omega$ by \noderef{MemW1pOn} and
$B_{R'}\subseteq\Omega$; $|g_v|^{p}$ is integrable on $B_{R'}$ because
$v\in W^{1,p}(B_{R'})$ has weak gradient $g_v$ (\noderef{MemW1pOn}), which is a
hypothesis of this node; the constant $c$ lies in $L^p$ of the bounded set
$B_\rho$; and hence $|g_v-c|^{p}$ is integrable on $B_\rho$ and $|g-g_v|^{p}$ on
$B_{R'}$, by (E1) applied to the two summands in each case.  Inserting
(4.16) of the thesis at $(\rho,R')$ --- legitimate by Step 1 --- and
(4.17) of the thesis at $R'$,
\[
\int_{B_\rho}|g-c|^{p}dx
\ \le\ 2^{p-1}\Bigl[c_3\Bigl(\frac\rho{R'}\Bigr)^{p\alpha}\rho^{\,n}
+R'^{\,n+\delta'}\Bigr]\Lambda ,
\]
and then (4.15) of the thesis at $R'$, i.e.\ $\Lambda\le c_1R'^{-\tau}$, the bracket
being nonnegative and $\Lambda\ge0$,
\[
\int_{B_\rho}|g-c|^{p}dx
\ \le\ 2^{p-1}c_1\Bigl[c_3\Bigl(\frac\rho{R'}\Bigr)^{p\alpha}\rho^{\,n}R'^{-\tau}
+R'^{\,n+\delta'-\tau}\Bigr],
\tag{4.18}
\]
which is the thesis' (4.18) of the thesis.

\emph{Step 3: the substitution $\rho=\frac12R'^{\,1+\theta}$.}  By Step 1,
\[
\frac\rho{R'}=\tfrac12R'^{\,\theta},
\qquad
\Bigl(\frac\rho{R'}\Bigr)^{p\alpha}=2^{-p\alpha}R'^{\,p\alpha\theta},
\qquad
\rho^{\,n}=2^{-n}R'^{\,n(1+\theta)} ,
\]
so the first term in the bracket of (4.18) equals
$c_32^{-p\alpha-n}R'^{\,p\alpha\theta+n+n\theta-\tau}
=c_32^{-p\alpha-n}R'^{\,n+n\theta+\eta}$, using $\eta=p\alpha\theta-\tau$; and
the second equals $R'^{\,n+\delta'-\tau}=R'^{\,n+n\theta+p\alpha\theta-\tau}
=R'^{\,n+n\theta+\eta}$, using $\delta'=n\theta+p\alpha\theta$.  Hence
\[
\int_{B_\rho}|g-c|^{p}dx
\ \le\ 2^{p-1}c_1\bigl(c_32^{-p\alpha-n}+1\bigr)\,R'^{\,n+n\theta+\eta} .
\]
Now $R'^{\,1+\theta}=2\rho$ gives
$R'^{\,n+n\theta}=\bigl(R'^{\,1+\theta}\bigr)^{n}=(2\rho)^{n}$ and
$R'^{\,\eta}=\bigl((2\rho)^{\frac1{1+\theta}}\bigr)^{\eta}
=(2\rho)^{\frac\eta{1+\theta}}=(2\rho)^{\lambda_0}$, so
\[
R'^{\,n+n\theta+\eta}=(2\rho)^{\,n+\lambda_0}
=2^{\,n+\lambda_0}\rho^{\,n+\lambda_0},
\]
an equality, not an estimate: because $\lambda_0$ is defined to be
$\eta/(1+\theta)$ exactly and is not truncated at $p$, no appeal to $\rho\le1$
is needed here, and none is made.  Combining,
\[
\int_{B_\rho}|g-c|^{p}dx
\ \le\ 2^{p-1}c_1\bigl(c_32^{-p\alpha-n}+1\bigr)2^{\,n+\lambda_0}
\,\rho^{\,n+\lambda_0}
\ =\ 2^{-p}c_{10}\,\rho^{\,n+\lambda_0},
\]
by the definition of $c_{10}$ in the preamble to this proof.  This is the thesis'
chain
$\int|g-c|^p\le c_8R'^{\,n+n\theta+\eta}\le c_9\rho^{\,n}\rho^{\eta/(1+\theta)}
\le c_{10}\rho^{\,n+\lambda_0}$.

\emph{Step 4: from the mean of $g_v$ to the mean of $\nabla u$.}  The ball
$B_\rho$ is measurable with $0<|B_\rho|=\omega_n\rho^{\,n}<\infty$, the field
$g$ belongs to $L^p(B_\rho)$ because $|g|^p$ is integrable on $\Omega$
(\noderef{MemW1pOn}) and $B_\rho\subseteq B_{R'}\subseteq\Omega$, and
$c\in\R^n$; so \noderef{MeanReplacementLp} applies with $S=B_\rho$, $f=g$ and
this $c$, and gives
\[
\int_{B_\rho(x_0)}\bigl|\nabla u-(\nabla u)_{B_\rho(x_0)}\bigr|^{p}dx
\ \le\ 2^{p}\int_{B_\rho}|g-c|^{p}dx .
\]
With Step 3,
\[
\int_{B_\rho(x_0)}\bigl|\nabla u-(\nabla u)_{B_\rho(x_0)}\bigr|^{p}dx
\ \le\ 2^{p}\cdot2^{-p}c_{10}\,\rho^{\,n+\lambda_0}
\ =\ c_{10}\,\rho^{\,n+\lambda_0}.
\]
This is (4.19) of the thesis.  Since $x_0$, $R$ and $\rho$ were arbitrary subject to
the stated constraints, the assertion follows.  (The hypothesis
$0\le\varepsilon\le1$ is not needed in this node's argument, and is retained
only because the bundle it is applied to is the one the theorems of \S4 and \S5
produce under that normalisation.)
\end{proof}
